\selectlanguage{ngerman}
\vspace*{-2em}
\par\noindent 
Die Reihe VIII der Leibniz-Edition ist ein durch das Akademienprogramm gefördertes Langzeitvorhaben der Berlin-Brandenburgischen Akademie der Wissenschaften. An der für sie eigens gegründeten Arbeitsstelle wird die Reihe VIII seit 2001 bearbeitetet. Sie legt hiermit ihren vierten Band vor, den letzten, der in Berlin seinen Abschluss findet, bevor alle weiteren ab 2025 in Potsdam entstehen werden. %
Die Reihe VIII ist derzeit auf zwölf Bände geplant. Dabei bilden nur die ersten beiden Bände der Reihe eine chronologische Einheit mit Schriften aus den Jahren, die Leibniz noch in Mainz, dann lange in Paris und auf Reisen nach London und Amsterdam verbracht hat (1672--1676). Alle späteren Bände dagegen folgen wie auch der vorliegende einer thematischen Ordnung und vereinen Schriften, die aus der gesamten Schaffenszeit von Leibniz stammen und nach Thema getrennt jeweils in chronologischer Reihenfolge ediert werden.%
\par%
Wir danken dem Bundesministerium für Bildung und Forschung sowie dem Regierenden Bürgermeister von Berlin, Senatskanzlei für Wissenschaft und Forschung, für die Finanzierung des Vorhabens. 
Arbeitsgrundlage der Berliner Editionsstelle sind die Digitalisate der in Reihe VIII zu edierenden Handschriften, die in hochauflösender Qualität und mit Durchlichtaufnahmen zur Erkennung von Wasserzeichen zum Start der Reihe eigens angefertigt worden waren. Wir verdanken diese digitale Arbeitsgrundlage der umfassenden Finanzierung seitens der Deutschen Forschungsgemeinschaft, der Alfried Krupp von Bohlen und Halbach-Stiftung und der Stiftung der Versicherungsgruppe Hannover sowie der freundlichen Unterstützung durch die Gottfried Wilhelm Leibniz Bibliothek Hannover (Niedersächsische Landesbibliothek). 
\par%
Auch für die Bearbeitung dieses vierten Bandes der Reihe bedanken wir uns wieder ganz herzlich bei der Gottfried Wilhelm Leibniz Bibliothek Hannover, deren Hilfe wir vielfach in Anspruch nehmen durften, insbesondere durch den ehemaligen Leiter der Abteilung Handschriften, Matthias Wehry, sowie durch die langjährige, leitende Mitarbeiterin Anja Fleck, die uns mit Auskünften zu den Beständen, der Durchsicht von Handschriften und deren Bereitstellung von Ort diente; bei Christoph Valentin, Abteilung Digitalisierung, bedanken wir uns für die ungezählten Digitalisate von Leibniz-Handschriften, mit denen er uns versorgt hat.
Vonseiten der Leibniz-Forschungsstelle Hannover haben wir vielfache Unterstützung erfahren: Charlotte Wahl, Siegmund Probst und Achim Trunk verdanken wir Hinweise auf weitere Handschriften für Reihe VIII sowie Informationen zu bereits edierten Stücken, die wir in den "Ergänzungen und Berichtigungen" dieses Bandes berücksichtigt haben; besonders in Siegmund Probst hatten wir einen verlässlichen Ansprechpartner in allen kniffligen Fragen, dessen beeindruckende Kompetenz wir mehrfach in Anspruch nehmen durften. Wir danken Stephan Waldhoff von der Editionsstelle Potsdam für seine Hilfe in Fragen der Datierung und bei der Identifizierung von Wasserzeichen und Schreiberhänden. Wir danken Stephan Meier-Oeser von der Editionsstelle Münster für die Retrodigitalisierung älterer Bände der Leibniz-Edition und deren Bereitstellung als PDF-Dateien sowie für eine Aktualisierung der aus den PDF-Dateien der Edition erzeugten txt-Dateien, mit denen eine Volltextsuche in den bislang edierten Stücken über Reihen und Bände hinweg möglich ist.
Wir danken Anne Steiger, unserer studentischen Mitarbeiterin, für ihre Unterstützung in verschiedenen Bereichen der Edition.
Unserer ehemaligen studentischen Mitarbeiterin, Pauline Just, sind wir weiterhin zu Dank verpflichtet: Auch für diesen Band war die von ihr erstellte Dokumentation von Wasserzeichen ein unerlässliches Werkzeug für die Datierung von Handschriften.
Gunthild Peters (geb. Storeck) danken wir für die im Rahmen ihrer Werkverträge geleisteten Vorarbeiten, die sich für die Bearbeitung von Stücken dieses Bandes nutzen ließen (N.\,\ref{49409}, N.\,\ref{41199}, N.\,\ref{58940}, N.\,\ref{58939}).
\par%
Außerhalb der Edition geht unser Dank an Richard Arthur und Osvaldo Ottaviani, die großzügig ihre noch nicht veröffentlichten Forschungsergebnisse für die Bearbeitung von N.\,\ref{60080} mit uns geteilt haben.%
\footnote{\, G.\,W. Leibniz, "Deriving the laws of motion from abstract principles. Translation and Commentary", hrsg. v. R. Arthur u. O. Ottaviani, \textit{Leibniz Review} XXXIV, 2024 [2025] (im Erscheinen).}
Bei Osvaldo Ottaviani bedanken wir uns darüber hinaus ganz herzlich für seine exzellenten Transkriptionen, die er uns für N.\,\ref{60349-1}, N.\,\ref{60349-2} und N.\,\ref{42488} frei zur Verfügung gestellt hat. Gedankt sei der Biblioteka Narodowa (Warschau) und der Universitätsbibliothek Wien für ihre Auskünfte zu Wasserzeichen und Leibnizbeständen.%
\par%
Wir danken der \LaTeX-Community und insbesondere Maïeul Rouquette (reledmac) für die fortlaufende Verbesserung des Textsatzprogrammes -- auch nach Wünschen der Reihe VIII. Für die Erstellung des vorliegenden Bandes wurde \LaTeX\ in der TeX Live Distribution von 2021 benutzt; die Zeichnungen sind mit der gemeinfreien Software Inkscape entstanden; für Zwecke der Datenauswertung, -überprüfung und -korrektur wurde das Sprachmodell von Anthropic (ClaudeAI) genutzt. %
Wir danken Serena Pirrotta vom De Gruyter Verlag für die gute Zusammenarbeit.
\par%
\vspace{2em}
Berlin, im Dezember 2024\hspace{65mm}Harald Siebert
